This chapter covers the infrastructure that surrounds the machine learning
models: how experiments are tracked, how models are promoted to production,
how the system detects when retraining is needed, and how operational health
is monitored. Together, these components close the loop between model serving
and model improvement.

% ═══════════════════════════════════════════════════════════
\section{MLOps Lifecycle Overview}
\label{sec:mlops-lifecycle}
% ═══════════════════════════════════════════════════════════

A trained model that cannot be retrained, versioned, or monitored is a
liability rather than an asset. The MLOps layer transforms the models
described in Chapter~\ref{ch:training} from static artifacts into a managed
lifecycle with four stages:

\begin{enumerate}
    \item \textbf{Train} --- execute a training job that produces a candidate
          model with logged hyperparameters and metrics.
    \item \textbf{Promote} --- compare the candidate against the current
          production model and replace it only if performance is equal or
          better.
    \item \textbf{Serve} --- the streaming pipeline loads the promoted model
          from a shared volume mount, requiring no image rebuild.
    \item \textbf{Monitor} --- observe the score distribution in production
          and trigger a new training cycle when drift is detected or
          supervisor feedback indicates degradation.
\end{enumerate}

The lifecycle is largely automated. Training, promotion, serving, and
monitoring all run without human intervention. The single human touch point
is the supervisor confirming or rejecting alerts in the dashboard
(Section~\ref{sec:feedback-loop}), which feeds one of the three retraining
triggers. The promotion gate in particular has no manual approval step
between a candidate model and production deployment --- a deliberate scope
decision discussed in Section~\ref{sec:promotion-gate}.

Three independent triggers can initiate a new training cycle: a weekly
schedule, a drift alert from the monitoring stack, and a supervisor feedback
divergence signal. Sections~\ref{sec:retrain-triggers}
and~\ref{sec:feedback-loop} describe each path.


% ═══════════════════════════════════════════════════════════
\section{Experiment Tracking with MLflow}
\label{sec:mlflow}
% ═══════════════════════════════════════════════════════════

MLflow provides the experiment tracking layer. Every training job --- whether
triggered by schedule, drift, or feedback --- logs its artifacts to the same
MLflow server, creating a single lineage of all model versions.

\subsection{Server Configuration}

The MLflow server runs as a Docker container with a flat-file backend
(\texttt{mlruns/} mounted as a volume) rather than a database-backed
tracking store. For a single-user development deployment processing one
retraining job per week, the flat-file backend avoids the operational
overhead of a dedicated metadata database. A production deployment serving
multiple concurrent training jobs would migrate to a PostgreSQL-backed store.

\subsection{Training Job Structure}

Each execution of the training script (\texttt{src/training/train\_all.py})
produces two independent MLflow runs under the experiment
\texttt{amen-anomaly-detection}: one for the Autoencoder and one for the
LSTM. The runs are independent rather than nested because the two models
have different architectures, hyperparameters, and evaluation metrics ---
nesting them would conflate unrelated configuration spaces.

Each run logs:

\begin{itemize}
    \item \textbf{Parameters} --- architecture dimensions, learning rate,
          batch size, number of epochs, sequence length (LSTM only).
    \item \textbf{Metrics} --- AUC-ROC on the validation set, computed
          against the oracle anomaly labels from the synthetic generator.
    \item \textbf{Artifacts} --- the trained model file (\texttt{.pt}),
          the fitted scaler, and reference score distributions
          (\texttt{ae\_ref\_scores.npy}, \texttt{lstm\_ref\_scores.npy})
          used by the drift detection system
          (Section~\ref{sec:retrain-triggers}).
\end{itemize}

% ═══════════════════════════════════════════════════════════
\section{Model Promotion Gate}
\label{sec:promotion-gate}
% ═══════════════════════════════════════════════════════════

After a training job completes, the system must decide whether the newly
trained model should replace the one currently serving in production. This
decision is handled by an automated promotion gate that compares validation
AUC scores.

\subsection{Promotion Logic}

Each model type (Autoencoder and LSTM) is evaluated independently. The
training script reads a file \texttt{models/production\_metrics.json} that
stores the AUC of the current production model for each type. If the
candidate's AUC is greater than or equal to the production AUC, the
candidate is promoted: its \texttt{.pt} file, scaler, and reference score
distribution are written to the shared model volume, and
\texttt{production\_metrics.json} is updated with the new baseline.

The comparison uses $\geq$ rather than $>$ to allow retraining on updated
data to promote a model even when the AUC is numerically identical. In
practice, a model retrained on fresher data with the same AUC is preferable
to one trained on older data, because the feature distributions it has seen
are more representative of current traffic.

On the very first training run --- when no production model exists ---
\texttt{production\_metrics.json} is absent. The training script detects
this and bootstraps by promoting unconditionally, creating the baseline
file for all subsequent comparisons.

\subsection{Scope and Limitations}

This gate is sufficient for the current single-user, single-environment
deployment. It has two deliberate limitations worth noting:

\begin{enumerate}
    \item \textbf{No human approval step.} A model that passes the AUC
          threshold is promoted immediately. In a production banking
          environment, a compliance or risk officer would review model
          changes before deployment. Adding a manual approval stage ---
          for example, an Airflow task that pauses the DAG until a human
          confirms via the dashboard --- is a straightforward extension.
    \item \textbf{Single-metric comparison.} The gate evaluates AUC only.
          A more robust gate would include additional criteria: F1-score
          at the operating threshold, per-scenario recall to ensure no
          regression on specific anomaly types, and inference latency to
          guard against deploying a model that is accurate but too slow
          for real-time scoring.
\end{enumerate}


% ═══════════════════════════════════════════════════════════
\section{Orchestration with Apache Airflow}
\label{sec:airflow}
% ═══════════════════════════════════════════════════════════

Chapter~\ref{ch:streaming} described the nightly profile rebuild DAG
(\texttt{nightly\_profiles}) that maintains the behavioral profiles consumed
by the streaming pipeline. This section focuses on the second DAG ---
\texttt{weekly\_retrain} --- and the infrastructure decisions that apply to
both.

\subsection{Airflow Infrastructure}

Airflow runs with a \textbf{LocalExecutor}, which executes tasks in the
scheduler process rather than distributing them to remote workers. This is
sufficient for the current workload of two DAGs with low concurrency.
Migrating to \texttt{CeleryExecutor} is a configuration change that
requires no DAG modifications --- only the addition of worker containers
and a message broker.

The Airflow metadata tables share the application PostgreSQL instance to
conserve memory on the development machine. A dedicated
\texttt{Dockerfile.airflow} installs only the packages required for
orchestration (\texttt{requirements-airflow.txt}: \texttt{redis},
\texttt{numpy}, \texttt{pandas}, \texttt{psycopg2-binary},
\texttt{python-dotenv}, \texttt{apache-airflow-providers-docker}) ---
excluding PyTorch, FastAPI, and SHAP. This keeps the Airflow image
approximately 2\,GB smaller than the training image.

\subsection{Weekly Retrain DAG}

The \texttt{weekly\_retrain} DAG runs on a weekly cron schedule and contains
three tasks:

\begin{enumerate}
    \item \textbf{Validate data freshness} --- confirm that the
          \texttt{transactions} table has received new events since the
          last training run. If no new data exists, the DAG short-circuits
          to avoid wasting compute on an identical training set.
    \item \textbf{Retrain models} --- executed via a
          \texttt{DockerOperator} that launches the training image
          (\texttt{amen-anomaly-training:latest}) as a separate container.
    \item \textbf{Verify promotion} --- check that
          \texttt{production\_metrics.json} was updated, confirming that
          at least one model was promoted.
\end{enumerate}

\subsection{Training Isolation with DockerOperator}

The retraining task uses Airflow's \texttt{DockerOperator} rather than a
\texttt{BashOperator} or \texttt{PythonOperator} for two reasons:

\begin{itemize}
    \item \textbf{Dependency isolation.} The training job requires PyTorch,
          which is large and unnecessary in the Airflow scheduler. Running
          training in its own container keeps the Airflow image lightweight.
    \item \textbf{Reproducibility.} The training image is pinned as
          \texttt{amen-anomaly-training:latest} with
          \texttt{force\_pull=False}, ensuring that the same image is used
          for scheduled, drift-triggered, and feedback-triggered retraining.
          The image is rebuilt explicitly when the training code changes,
          not pulled automatically on every run.
\end{itemize}

The \texttt{DockerOperator} is configured with
\texttt{mount\_tmp\_dir=False} to prevent Airflow from creating a temporary
directory mount that conflicts with the container's working directory. The
model volume (\texttt{./models}) is mounted into the training container at
the same path used by the scoring service, so a promoted model is
immediately available to the streaming pipeline on its next restart ---
no file copy or deployment step required.
% ═══════════════════════════════════════════════════════════
\section{Retraining Triggers}
\label{sec:retrain-triggers}
% ═══════════════════════════════════════════════════════════

The system supports three independent paths to initiate a retraining cycle.
Each path ultimately triggers the same \texttt{weekly\_retrain} DAG
described in Section~\ref{sec:airflow}, ensuring that all retraining jobs
--- regardless of their origin --- follow the same training, promotion, and
artifact-logging pipeline.

\subsection{Scheduled Retraining}

The simplest trigger is a weekly cron schedule configured directly in the
Airflow DAG definition. This provides a baseline retraining cadence that
ensures the models are periodically refreshed with recent transaction data,
even in the absence of any detected anomaly.

\subsection{Drift-Triggered Retraining}

The second trigger detects \textit{score distribution drift} --- a
sustained shift in the anomaly scores produced by the scoring service that
suggests the models are no longer well-calibrated to the current data.

During each training job, the training script saves the distribution of
anomaly scores on the validation set as reference arrays
(\texttt{ae\_ref\_scores.npy}, \texttt{lstm\_ref\_scores.npy}), logged as
MLflow artifacts (Section~\ref{sec:mlflow}). In production, the scoring
service exposes the fraction of events scoring above the T1 threshold
($0.95$) as a Prometheus gauge. A Prometheus alert rule
(\texttt{config/alert\_rules.yml}) fires when this fraction exceeds 15\%
for a sustained period of one hour:

\begin{verbatim}
- alert: ScoreDistributionDrift
  expr: score_above_t1_ratio > 0.15
  for: 1h
\end{verbatim}

When this alert fires, the detection chain involves four components:

\begin{enumerate}
    \item \textbf{Prometheus} evaluates the alert rule and forwards the
          firing alert to AlertManager.
    \item \textbf{AlertManager} (port 9093) receives the alert, deduplicates
          it with a \texttt{repeat\_interval} of 24 hours to prevent
          repeated triggers from the same drift episode, and routes it to
          a webhook receiver.
    \item \textbf{Webhook relay} (\texttt{src/webhook/relay.py}) --- a
          lightweight FastAPI service that translates the AlertManager
          webhook payload into an Airflow REST API call. This intermediary
          exists because AlertManager's native webhook format does not
          support the HTTP Basic Authentication required by Airflow's REST
          API. The relay is a single-endpoint service
          (Dockerfile built on \texttt{python:3.11-slim}, approximately
          150\,MB) that performs authentication and forwards the trigger
          with \texttt{conf=\{"trigger": "drift\_alert"\}}.
    \item \textbf{Airflow REST API} receives the trigger and starts the
          \texttt{weekly\_retrain} DAG. The API is configured with
          \texttt{basic\_auth} as the authentication backend.
\end{enumerate}
\begin{figure}[H]
\centering
\begin{tikzpicture}[
    node distance=0.8cm,
    block/.style={draw, rectangle, rounded corners=3pt, fill=gray!8,
                  minimum width=3.5cm, minimum height=1cm,
                  font=\small, align=center, thick},
    arrow/.style={-{Stealth[length=3mm]}, thick},
    note/.style={font=\scriptsize\itshape, text width=3.2cm, align=center}
]

% Nodes
\node[block] (prom) {Prometheus\\{\scriptsize Alert rule evaluation}};
\node[block, below=of prom] (am) {AlertManager\\{\scriptsize Port 9093}};
\node[block, below=of am] (webhook) {Webhook Relay\\{\scriptsize FastAPI, port 5005}};
\node[block, below=of webhook] (airflow) {Airflow REST API\\{\scriptsize Basic auth}};
\node[block, below=of airflow] (dag) {\texttt{weekly\_retrain}\\{\scriptsize DAG execution}};

% Arrows with annotations
\draw[arrow] (prom) -- node[right, note] {Alert fires after\\1h above 15\%} (am);
\draw[arrow] (am) -- node[right, note] {Webhook POST\\24h dedup} (webhook);
\draw[arrow] (webhook) -- node[right, note] {Auth translation\\+ trigger conf} (airflow);
\draw[arrow] (airflow) -- node[right, note] {\texttt{conf=\{"trigger":\\"drift\_alert"\}}} (dag);

\end{tikzpicture}
\caption{Drift-triggered retraining chain. The webhook relay exists because
         AlertManager cannot natively authenticate against Airflow's REST
         API.}
\label{fig:drift-chain}
\end{figure}

The 15\% threshold and one-hour window were chosen empirically: under
normal operation with the current synthetic data, fewer than 5\% of events
exceed T1. A sustained elevation above 15\% indicates that the score
distribution has shifted materially, not that a single burst of legitimate
anomalies occurred.


\subsection{Feedback-Triggered Retraining}

The third trigger is described in Section~\ref{sec:feedback-loop}, as it
depends on the supervisor feedback mechanism.


% ═══════════════════════════════════════════════════════════
\section{Supervisor Feedback Loop}
\label{sec:feedback-loop}
% ═══════════════════════════════════════════════════════════

The feedback loop is the single human touch point in the MLOps lifecycle.
When the streaming pipeline produces an alert at the ALERT or BLOCK tier,
the branch supervisor reviews it in the dashboard
(Chapter~\ref{ch:dashboard}) and either confirms or rejects it. This
decision is submitted to the simulation API via a
\texttt{POST /api/v1/feedback} endpoint, which writes the supervisor's
verdict to the \texttt{evaluation\_labels} table in PostgreSQL with a
safe foreign key reference to the original transaction.

These labels serve two purposes. First, they provide ground truth for
evaluating model performance on real decisions --- a capability that
becomes meaningful once the system operates on real transaction data
rather than synthetic events with oracle labels. Second, they feed the
feedback divergence DAG.

\subsection{Feedback Divergence DAG}

An Airflow DAG (\texttt{dags/feedback\_monitor.py}) runs daily and executes
a SQL query against the \texttt{evaluation\_labels} table: over a rolling
seven-day window, what fraction of ALERT and BLOCK tier events were
rejected by supervisors? If the rejection rate exceeds 50\% with a minimum
of 10 reviewed events, the DAG triggers the \texttt{weekly\_retrain} DAG
with \texttt{conf=\{"trigger": "feedback\_divergence"\}}.

The 50\% threshold reflects a simple principle: if more than half of the
system's high-confidence alerts are being overridden by human reviewers,
the models are producing more noise than signal at the upper tiers and
should be retrained. The minimum of 10 reviews prevents a single rejected
alert early in a deployment from triggering an unnecessary retraining
cycle.

\subsection{Three Triggers, One Pipeline}

Table~\ref{tab:retrain-triggers} summarizes the three retraining paths.

\begin{table}[H]
\centering
\caption{Retraining trigger comparison}
\label{tab:retrain-triggers}
\begin{tabular}{@{}llll@{}}
\toprule
\textbf{Trigger} & \textbf{Source} & \textbf{Condition} & \textbf{Frequency} \\
\midrule
Scheduled    & Airflow cron        & Weekly              & Fixed \\
Drift        & Prometheus alert    & $>$15\% above T1 for 1h & On demand \\
Feedback     & Supervisor reviews  & $>$50\% rejection, $\geq$10 reviews & Daily check \\
\bottomrule
\end{tabular}
\end{table}

All three paths converge on the same \texttt{weekly\_retrain} DAG. The
\texttt{conf} parameter passed at trigger time records which path initiated
the run, providing traceability in the Airflow logs without requiring
separate DAG definitions for each trigger type.


% ═══════════════════════════════════════════════════════════
\section{Observability and Monitoring}
\label{sec:observability}
% ═══════════════════════════════════════════════════════════

Observability covers two complementary concerns: \textit{metrics} for
real-time dashboarding and alerting, and \textit{logs} for post-incident
diagnosis.

\subsection{Service Metrics}

Each streaming service (enrichment, scoring, decision) embeds a
\texttt{ServiceMetrics} class from \texttt{src/monitoring/metrics.py} that
exposes two Prometheus instruments:

\begin{itemize}
    \item A \textbf{Counter} tracking the total number of events processed,
          labeled by service name and outcome (success or error).
    \item A \textbf{Histogram} measuring per-event processing latency in
          seconds, enabling percentile queries (p50, p95, p99) in Grafana.
\end{itemize}

Each service exposes its metrics on a dedicated port: enrichment on 9100,
scoring on 9101, and decision on 9102. A Kafka exporter container (port
9308) provides broker-level metrics including consumer group lag --- the
difference between the latest produced offset and the latest committed
offset --- which serves as the primary indicator of pipeline backpressure.
Prometheus scrapes all four targets at a configurable interval
(\texttt{config/prometheus.yml}).

\subsection{Failure Domain Separation}

A monitoring system must not share a failure domain with the service it
monitors. If the enrichment service and its health reporter both depend on
Redis, a Redis outage would simultaneously break the service \textit{and}
suppress the alert that should report the breakage. This principle ruled
out Redis-based heartbeat mechanisms early in the design process.

The chosen architecture avoids this problem: Prometheus pulls metrics over
HTTP from each service's embedded metrics endpoint. The only shared
dependency is the network layer itself, which is an acceptable residual
risk --- if the network is down, no monitoring architecture can report it
from within the affected network.

\subsection{Structured Logging}

All streaming services use a \texttt{JsonFormatter} class
(\texttt{src/streaming/logger.py}) that outputs one JSON object per log
line. Each line includes a timestamp, service name, log level, and message,
with optional contextual fields such as \texttt{event\_id},
\texttt{client\_id}, and \texttt{error\_type} attached as structured
extras.

A \texttt{setup\_logger(service\_name)} factory configures Python's
standard \texttt{logging} module with the JSON formatter, replacing all
\texttt{print()} statements that existed in earlier versions of the
codebase. Structured JSON logs enable downstream processing --- filtering
by service, grouping by error type, or correlating events across services
by \texttt{event\_id} --- without fragile regular expression parsing of
human-readable log strings.